\documentclass[12pt,a4paper]{article}

\usepackage[utf8]{inputenc}
\usepackage[french]{babel}
\usepackage{graphicx}
\usepackage{amsmath}
\usepackage{geometry}
\geometry{margin=2.5cm}

\title{Analyse de la consommation énergétique des Centres de Données en Mauritanie}
\author{Rokia  \And Sara \And Aicha }

\date{\today}

\begin{document}

\maketitle

\begin{abstract}
Les centres de données constituent l’infrastructure essentielle du numérique moderne. 
En Mauritanie, le développement des télécommunications et des services en ligne 
entraîne une augmentation progressive de la demande énergétique. 
Ce travail propose une analyse comparative de la consommation énergétique 
de plusieurs centres de données opérant localement. 
L’objectif est d’identifier les facteurs influençant la consommation, 
notamment la capacité de stockage, le nombre de serveurs et le type de refroidissement, 
afin de proposer des pistes d’amélioration de l’efficacité énergétique.
\end{abstract}

\section{Introduction}

La transformation numérique a conduit à une croissance rapide des centres de données. 
Ces infrastructures consomment une quantité importante d’électricité 
pour alimenter les serveurs et maintenir une température optimale. 

En Mauritanie, l’expansion des opérateurs télécoms et des services cloud 
rend nécessaire l’analyse de leur impact énergétique.

\section{Problématique}

Comment analyser la consommation énergétique des centres de données en Mauritanie 
et identifier les leviers d’optimisation tout en maintenant la performance des services ?

\section{Méthodologie}

Cette étude repose sur une analyse comparative basée sur des estimations académiques. 
Les données suivantes sont utilisées à des fins pédagogiques.

L’indicateur principal étudié est le PUE (Power Usage Effectiveness) :

\[
PUE = \frac{\text{Énergie totale consommée}}{\text{Énergie utilisée par les équipements IT}}
\]


\begin{figure}[h]
\centering
\includegraphics[width=1\textwidth]{data.png}
\caption{Comparaison de la consommation énergétique des centres de données en Mauritanie}
\end{figure}



\section{Résultats et Analyse}

Les centres disposant d’un système de refroidissement liquide 
présentent généralement une meilleure efficacité énergétique. 
Bien que leur consommation totale soit plus élevée, 
leur performance par unité de stockage est plus optimisée.

Les centres utilisant un refroidissement à air 
affichent une consommation proportionnellement plus importante 
par rapport à leur capacité.

\begin{figure}[h]
\centering
\includegraphics[width=0.8\textwidth]{diagramme.png}
\caption{Diagramme comparatif de la consommation énergétique des centres de données en Mauritanie}
\end{figure}

Le diagramme met en évidence les différences de consommation 
entre les centres étudiés. 
On observe que les centres de grande capacité consomment davantage d’énergie, 
mais bénéficient d’une meilleure optimisation technologique.

\section{Conclusion}

L’analyse montre que la consommation énergétique des centres de données 
dépend principalement de la capacité installée et du système de refroidissement. 

L’amélioration de l’efficacité énergétique peut passer par :
\begin{itemize}
\item L’adoption du refroidissement liquide,
\item L’optimisation de la gestion thermique,
\item L’intégration d’énergies renouvelables.
\end{itemize}

Cette étude met en évidence l’importance d’une gestion énergétique responsable 
dans le développement numérique en Mauritanie.

\begin{thebibliography}{9}

\bibitem{ref1}
International Energy Agency, Data Centres and Energy Report, 2023.

\bibitem{ref2}
A. Shehabi et al., Data Center Energy Usage Report, 2022.

\end{thebibliography}

\end{document}
